\subsection*{11.5 Strong Law of Large Numbers}

If we assume that the fourth moments of the random variables are bounded, we have

an easy proof of the strong law of large number.

\begin{thmbox}{11.7 (4th-moment Strong Law of Large Numbers)}
\end{thmbox}

Suppose Xi ,f o r i\geq 1 , are independent random variables with mean \mu and

E[X4

i ]\leq c< \infty for all i Then Sn

n \to \mu almost surely.

\textit{Proof} Without loss of generality, we assume that \mu= 0 . (We can consider Yi =

Xi -\mu if the mean of Xi is not zero.)

The expectation of S4

n can be expanded as

E[S4

n]= E[(X1 +X 2 +\cdot\cdot\cdot+ Xn)4]

=E

[ n\sum

i=1

X4

i +3

\sum

i/=j

X2

i X2

j +4

\sum

i/=j

XiX3

j

+6

\sum

i,j,k

i,j,kdistinct

XiXj X2

k +

\sum

i,j,k,\ell

i,j,k,\elldistinct

XiXj XkX\ell

]

We analyze the terms one by one. For distinct indices i , j , k , and \ell,we have

E[XiXj XkX\ell]= E[XiXj X2

k]= E[XiX3

j ]= 0,

by the assumption that X1,...,X n are independent. The fourth-power terms are

bounded byE[\sumn

i=1 X4

i ]\leq nc becauseE[X4

i ]\leq c.

Fori/=j , by Cauchy-Schwarz inequality,

E[X2

i X2

j ]\leq

\sqrt

E[X4

i ]E[X4

j ]\leq \sqrt c\cdot c = c.

Hence,

3E

[ \sum

i/=j

X2

i X2

j

]

\leq3 n(n- 1 )c

This givesE[S4

n]\leq nc+ 3 n(n- 1 )c. We now fix any \epsilon> 0,

P

( \vertSn\vert

n >\epsilon

)

=P(S 4

n >n 4\epsilon4)\leq nc+ 3 n(n- 1 )c

n4\epsilon4 \leq 3n2c+ nc - 3 nc

n4\epsilon4 \leq 3c

n2\epsilon4 .

Because

\infty\sum

n=1

P

( \vertSn\vert

n >\epsilon

)

\leq

\infty\sum

n=1

3c

n2\epsilon4 <\infty ,

we can apply the first Borel-Cantelli lemma (Theorem 5.8) to conclude that the

event .{\vertSn\vert/n > \epsilon io. } has probability 0. Therefore, by Theorem 10.1, Sn/n

converges to 0 almost surely. \hfill $\square$

Note that Theorem 11.7 does not assume that the random variables Xi are

identically distributed, but it requires that the 4th moments are uniformly bounded.

We state below a stronger version the strong law that assumes finite mean and

pairwise independence.

\begin{thmbox}{11.8}
\end{thmbox}

Suppose Xi ,f o r i\geq 1 , are pairwise independent and identically distributed

random variables with finite mean \mu. ThenSn/n

as.

\to \mu asn\to\infty .

A proof of this theorem by Etemadi can be found in [ 4, Thm 2.4.1].

\textbf{Example 11.5.1 (Estimation of Cumulative Distribution Function)} \\

Consider the problem of estimating the cdf F(x) = Pr (X\leq x) of a random variable X by

generating iid. samples Xn from the cdf F for n\geq 1 We can estimate F(x) by the empirical

distribution

Fn(x)\coloneqq 1

n

n\sum

k=1

1[Xk,\infty )(x).

The indicator function .1[Xk,\infty )(x) is equal to 1 if x\geq X k and is 0 otherwise. The function Fn first

counts the number of samples that is less than or equal to x , and then divides it by n Because it

depends on the random samples, the estimate Fn(x) is a random variable. By the strong law of

large numbers,Fn(x) converges to the correct value F(x) with probability 1.

This result is strengthened in the Glivenko-Cantelli theorem, which states that the convergence is

uniform in x ,

sup

x\inR

\vert Fn(x)- F(x) \vert

as.

-\to 0 .

The Glivenko-Cantelli theorem is informally known as "the fundamental theorem of statistics".

\textbf{Example 11.5.2 (Normal Numbers)} \\

A real number is said to be normal in base 10 if the frequency of each digit 0 through 9 in its

decimal expansion is asymptotically equal. If a number has multiple decimal expansions, such as

0.1 = 0.09999... , we choose the one that does not end in infinitely many trailing 0's. The strong

law of large numbers implies that, with probability 1, a randomly chosen number in the interval

[0, 1] is normal in base 10.

We can generalize this notion of normality to other bases. A real number is said to be normal if

it is normal in base b all for all integer b \geq 2 simultaneously. By the strong law of large numbers,

for any integer b \geq 2, there exists a set Ab in .[0, 1] with P(Ab) = 1 such that all numbers in Ab

are normal in base b. Taking the intersection of countably events with probability 1, we deduce

that a number in .[0, 1] is normal almost surely. This proves the existence of normal numbers.

However, the argument in the previous paragraph is non-constructive. It is a non-trivial task to

explicitly construct a normal number or prove that a given number is normal.
